\chapter{General Formalism of the Effective Action}
\label{ch:eff-action-general}

The central object of this thesis is the \emph{effective action} $\Gamma[\varphi_c]$,
a functional that encodes the full quantum dynamics of a field theory in a single
object. Before defining it, we build up to it systematically: starting from the
generating functional $Z[J]$ for all Green functions, passing through the
connected generating functional $W[J]$, and finally arriving at $\Gamma[\varphi_c]$
via a Legendre transform. The payoff is that the equation of motion derived from
$\Gamma$ governs the quantum-corrected vacuum structure of the theory, making it
the natural tool for studying spontaneous symmetry breaking and phase transitions.


\section{Generating Functionals}
\label{sec:generating-functionals}

Rather than computing Green functions diagram by diagram, all of them, full,
connected, and one-particle-irreducible, can be extracted from a small number
of functionals by simple functional differentiation.

Coupling the field $\varphi(x)$ linearly to an arbitrary $c$-number source $J(x)$,
\begin{equation}
    \mathcal{L}(\varphi,\partial_\mu\varphi)
    \;\longrightarrow\;
    \mathcal{L} + J(x)\,\varphi(x),
    \label{eq:source-coupling}
\end{equation}
the generating functional for the full Green functions is defined as the
vacuum-to-vacuum transition amplitude in the presence of this source,
\begin{equation}
    Z[J] = \langle 0 | T \exp\!\left( i \int J(x)\, \varphi(x)\, d^4x \right) | 0 \rangle,
    \label{eq:Z-operator-def}
\end{equation}
where $|0\rangle$ is the physical vacuum and $\varphi(x)$ is the Heisenberg-picture
field, both defined at $J=0$. Expanding in powers of $J$ makes the content of
this definition explicit,
\begin{align}
    Z[J] &= \sum_{n=0}^{\infty} \frac{i^n}{n!}
             \int d^4x_1 \cdots d^4x_n\;
             J(x_1) \cdots J(x_n)\;
             G_n(x_1, \ldots, x_n),
    \label{eq:Z-power-series}
\end{align}
where $G_n(x_1,\ldots,x_n) \equiv \langle 0|T[\varphi(x_1)\cdots\varphi(x_n)]|0\rangle$
is the full $n$-point Green function. Each $G_n$ is recovered by $n$ functional
derivatives of $Z[J]$ with respect to $J$, evaluated at $J=0$. Equivalently,
$Z[J]$ admits a path-integral representation,
\begin{equation}
    Z[J] = N \int [\mathcal{D}\varphi]\;
           \exp\!\left( i \int d^4x\,
           \bigl[\mathcal{L}[\varphi(x)] + J(x)\,\varphi(x)\bigr] \right),
    \label{eq:Z-path-integral}
\end{equation}
where $N$ is fixed by $Z[0]=1$. The two representations
\eqref{eq:Z-operator-def} and \eqref{eq:Z-path-integral} describe the same
object and will be used interchangeably throughout.

The full Green functions $G_n$ contain disconnected contributions, products
of lower-point functions, that carry no new dynamical information. The
\emph{connected} Green functions $G_n^c$, which isolate only the irreducible
correlations, are generated by passing to the logarithm of $Z[J]$,
\begin{equation}
    iW[J] = \ln Z[J],
    \label{eq:W-from-Z}
\end{equation}
the proof that this logarithm precisely removes all disconnected pieces is
given in Appendix~\ref{app:connected}. Since $Z[J]$ is the vacuum-to-vacuum
amplitude, Eq.~\eqref{eq:W-from-Z} means
\begin{equation}
    e^{iW[J]} = \langle 0^+|0^-\rangle_J,
    \label{eq:W-vacuum-persistence}
\end{equation}
where $|0^\pm\rangle$ are the vacuum states in the far future and far past, so
$W[J]$ is the \emph{vacuum persistence amplitude} in the presence of the
external source. This interpretation will be central when we discuss the
metastability of the false vacuum in Chapter~\ref{ch:vacuum-decay}.

The first functional derivative of $W[J]$ defines the \emph{classical field},
\begin{equation}
    \varphi_c(x) \equiv \frac{\delta W}{\delta J(x)}
    = \left[\frac{\langle 0^+|\varphi(x)|0^-\rangle}
                 {\langle 0^+|0^-\rangle}\right]_J,
    \label{eq:classical-field-def}
\end{equation}
the source-dependent vacuum expectation value of $\varphi(x)$, which will
serve as the argument of the effective action.


\section{The Effective Action and the Effective Potential}
\label{sec:effective-action}

Given $W[J]$ and the classical field $\varphi_c$ via \eqref{eq:classical-field-def},
we trade the source $J$ for $\varphi_c$ as the fundamental variable via the
Legendre transform,
\begin{equation}
    \Gamma[\varphi_c] = W[J] - \int d^4x\, J(x)\,\varphi_c(x),
    \label{eq:effective-action-def}
\end{equation}
where $J$ on the right-hand side is understood as a functional of $\varphi_c$
obtained by inverting \eqref{eq:classical-field-def}. Differentiating with
respect to $\varphi_c$ and using \eqref{eq:classical-field-def} gives
\begin{equation}
    \frac{\delta\Gamma}{\delta\varphi_c(x)} = -J(x).
    \label{eq:gamma-eom}
\end{equation}
In the physical situation $J=0$, this becomes the quantum equation of motion,
\begin{equation}
    \frac{\delta\Gamma}{\delta\varphi_c(x)} = 0,
    \label{eq:gamma-eom-J0}
\end{equation}
the exact analogue of $\delta S/\delta\varphi = 0$ for the classical action,
but now with all quantum corrections included. The vacua of the quantum theory
are precisely the solutions to \eqref{eq:gamma-eom-J0}.

Expanding $\Gamma[\varphi_c]$ in powers of $\varphi_c$,
\begin{equation}
    \Gamma[\varphi_c] = \sum_{n=0}^{\infty} \frac{1}{n!}
    \int d^4x_1\cdots d^4x_n\;
    \Gamma^{(n)}(x_1,\ldots,x_n)\;
    \varphi_c(x_1)\cdots\varphi_c(x_n),
    \label{eq:gamma-1PI-expansion}
\end{equation}
the coefficients $\Gamma^{(n)}$ are the one-particle-irreducible (1PI) Green
functions, or proper vertices, the sum of all Feynman diagrams with $n$ external
lines that cannot be disconnected by cutting a single internal propagator, with
no external propagators attached.

For a translationally invariant configuration $\varphi_c = \text{const}$, the
effective action takes the form $\Gamma = -\Omega\,V(\varphi_c)$, where $\Omega$
is the four-volume. More generally, for slowly-varying fields, $\Gamma$ admits
a derivative expansion,
\begin{equation}
    \Gamma[\varphi_c] = \int d^4x\left[
        -V(\varphi_c)
        + \frac{1}{2}(\partial_\mu\varphi_c)^2\, Z(\varphi_c)
        + \cdots\right],
    \label{eq:gamma-derivative-expansion}
\end{equation}
where $V(\varphi_c)$ is the \emph{effective potential}, an ordinary function
rather than a functional. Comparing with \eqref{eq:gamma-1PI-expansion}, the
$n$-th derivative of $V$ equals the sum of all 1PI diagrams with $n$ external
legs at zero external momenta, so the effective potential encodes the full
quantum-corrected vacuum structure of the theory. The physical mass and coupling
are read off from it via
\begin{equation}
    \mu^2 = \frac{d^2V}{d\varphi_c^2}\bigg|_{\varphi_c=0},
    \qquad
    \lambda = \frac{d^4V}{d\varphi_c^4}\bigg|_{\varphi_c=0}.
    \label{eq:renorm-conditions}
\end{equation}


\section{Spontaneous Symmetry Breaking}
\label{sec:ssb}

The effective potential is the natural language for spontaneous symmetry
breaking. Suppose $\mathcal{L}$ has a discrete internal symmetry,
e.g.\ $\varphi \to -\varphi$. This symmetry is spontaneously broken when the
quantum equation of motion \eqref{eq:gamma-eom-J0} admits a solution with a
nonzero, spatially uniform vacuum expectation value,
\begin{equation}
    \frac{dV}{d\varphi_c}\bigg|_{\varphi_c = \langle\varphi\rangle} = 0,
    \qquad \langle\varphi\rangle \neq 0,
    \label{eq:ssb-condition-potential}
\end{equation}
even in the absence of any external source. Stability requires this stationary
point to be a minimum of $V$.

Two phases are possible. In the symmetric phase, $\langle\varphi\rangle = 0$
and the $\varphi\to-\varphi$ symmetry is intact. In the broken phase,
$\langle\varphi\rangle = v \neq 0$ and the symmetry is hidden. The physical
spectrum in the broken phase is described by the shifted field
$\varphi' = \varphi - v$, and the renormalization conditions
\eqref{eq:renorm-conditions} are displaced to the broken minimum,
\begin{equation}
    \mu^2 = \frac{d^2V}{d\varphi_c^2}\bigg|_{\varphi_c=v},
    \qquad
    \lambda = \frac{d^4V}{d\varphi_c^4}\bigg|_{\varphi_c=v}.
    \label{eq:ssb-mass-coupling}
\end{equation}
With the formalism in place, we turn in the next chapter to an explicit
one-loop computation of the effective potential, first in a simple scalar
theory and then in the physically relevant case of scalar electrodynamics.
